\section{Leptonium production}
\label{sec:leptionium}

The QED analogues of quarkonia are bound states composed of a charged lepton and an antilepton. The six possible flavour combinations of purely leptonic atoms, collectively referred to as leptonia, are positronium ($e^-e^+$), muonium ($e^{-}\mu^{+}$), tauonium ($e^{-}\tau^{+}$), dimuonium ($\mu^-\mu^+$), mu--tauonium ($\mu^{-}\tau^{+}$), and ditauonium ($\tau^-\tau^+$), together with their charge-conjugated counterparts where applicable. Unlike NRQCD, the LDMEs governing bound-state formation in NRQED can be computed perturbatively. The explicit expressions for the \swave- and {\pwavetxt} leptonium LDMEs, derived from the Coulomb potential, are given in eqs.~\eqref{eq:ldme_leptonia_swave} and \eqref{eq:ldme_leptonia_pwave} in appendix~\ref{app:setup}, respectively.

%%%%%%%%%%%%%%%%%%%%%%%%%%%%%%%%%%%%%%%%%%%%%%%%%%%%%%%%%%%%%%%%
\begin{figure}[t]
    \centering
    \begin{subfigure}[b]{\textwidth}
        \centering
        \def\svgwidth{0.4\textwidth}
        % Created with FeynGame
% Include with input{lep_1.tex}
% Width can be set with \def\svgwidth{...} and a scaling with \def\svgscale{...}
% These options have to be set before the import statement and only apply to one included file
\begingroup%
  \makeatletter%
  \providecommand\rotatebox[2]{#2}%
  \ifx\svgwidth\undefined%
    \setlength{\unitlength}{595.32000732bp}%
    \ifx\svgscale\undefined%
      \relax%
    \else%
      \setlength{\unitlength}{\unitlength * \real{\svgscale}}%
    \fi%
  \else%
    \setlength{\unitlength}{\svgwidth}%
  \fi%
  \global\let\svgwidth\undefined%
  \global\let\svgscale\undefined%
  \makeatother%
  \parbox[c][][c]{\unitlength}{%
    \begin{picture}(1,0.6082474226804123)%
      \put(0,0){\includegraphics[width=\unitlength,viewport = 0 0 213 129]{feynman/lep_1.eps}}%
      \put(0.1030188665107254,0.05773196146660244){\raisebox{-0.5\height}{\rotatebox{0.0}{\scalebox{1.0}{\makebox(0,0)[]{{$e^+$}}}}}}%
      \put(0.7055781184237487,0.07712815942100788){\raisebox{-0.5\height}{\rotatebox{0.0}{\scalebox{1.0}{\makebox(0,0)[]{{$e^+$}}}}}}%
      \put(0.7657582486659001,0.41258839846176576){\raisebox{-0.5\height}{\rotatebox{0.0}{\scalebox{1.0}{\makebox(0,0)[]{{$e^-$}}}}}}%
      \put(0.5695876288659794,0.3015463917525773){\raisebox{-0.5\height}{\rotatebox{0.0}{\scalebox{1.0}{\makebox(0,0)[]{{$\gamma^\ast$}}}}}}%
      \put(0.1030188665107254,0.5703504050195699){\raisebox{-0.5\height}{\rotatebox{0.0}{\scalebox{1.0}{\makebox(0,0)[]{{$e^-$}}}}}}%
      \put(0.8776280436554367,0.564420488002808){\raisebox{-0.5\height}{\rotatebox{0.0}{\scalebox{1.0}{\makebox(0,0)[]{{$\gamma$}}}}}}%
      \put(0.9177938144329897,0.18814432989690721){\raisebox{-0.5\height}{\rotatebox{0.0}{\scalebox{1.0}{\makebox(0,0)[]{{$\dfrac{h_e}{\chi_{e J}}$}}}}}}%
    \end{picture}%
  }%
\endgroup%

        \def\svgwidth{0.4\textwidth}
        % Created with FeynGame
% Include with input{lep_2.tex}
% Width can be set with \def\svgwidth{...} and a scaling with \def\svgscale{...}
% These options have to be set before the import statement and only apply to one included file
\begingroup%
  \makeatletter%
  \providecommand\rotatebox[2]{#2}%
  \ifx\svgwidth\undefined%
    \setlength{\unitlength}{595.32000732bp}%
    \ifx\svgscale\undefined%
      \relax%
    \else%
      \setlength{\unitlength}{\unitlength * \real{\svgscale}}%
    \fi%
  \else%
    \setlength{\unitlength}{\svgwidth}%
  \fi%
  \global\let\svgwidth\undefined%
  \global\let\svgscale\undefined%
  \makeatother%
  \parbox[c][][c]{\unitlength}{%
    \begin{picture}(1,0.6223404255319149)%
      \put(0,0){\includegraphics[width=\unitlength,viewport = 0 0 213 132]{feynman/lep_2.eps}}%
      \put(0.7582824480914075,0.4257561133062902){\raisebox{-0.5\height}{\rotatebox{0.0}{\scalebox{1.0}{\makebox(0,0)[]{{$e^-$}}}}}}%
      \put(0.0930188665107254,0.5703504050195699){\raisebox{-0.5\height}{\rotatebox{0.0}{\scalebox{1.0}{\makebox(0,0)[]{{$e^-$}}}}}}%
      \put(0.08111702825160737,0.05957447087511103){\raisebox{-0.5\height}{\rotatebox{0.0}{\scalebox{1.0}{\makebox(0,0)[]{{$e^+$}}}}}}%
      \put(0.607714529690966,0.07409469690687043){\raisebox{-0.5\height}{\rotatebox{0.0}{\scalebox{1.0}{\makebox(0,0)[]{{$e^+$}}}}}}%
      \put(0.35106382978723405,0.31117021276595747){\raisebox{-0.5\height}{\rotatebox{0.0}{\scalebox{1.0}{\makebox(0,0)[]{{$\gamma^\ast$}}}}}}%
      \put(0.8776280436554367,0.564420488002808){\raisebox{-0.5\height}{\rotatebox{0.0}{\scalebox{1.0}{\makebox(0,0)[]{{$\gamma$}}}}}}%
      \put(0.9179148936170213,0.19414893617021275){\raisebox{-0.5\height}{\rotatebox{0.0}{\scalebox{1.0}{\makebox(0,0)[]{{$\dfrac{h_e}{\chi_{e J}}$}}}}}}%
    \end{picture}%
  }%
\endgroup%

        \caption{}\label{fig:diag_leptonium_a}
    \end{subfigure}
    \begin{subfigure}[b]{0.4\textwidth}
        \centering
        \def\svgwidth{\textwidth}
        % Created with FeynGame
% Include with input{lep_4.tex}
% Width can be set with \def\svgwidth{...} and a scaling with \def\svgscale{...}
% These options have to be set before the import statement and only apply to one included file
\begingroup%
  \makeatletter%
  \providecommand\rotatebox[2]{#2}%
  \ifx\svgwidth\undefined%
    \setlength{\unitlength}{595.32000732bp}%
    \ifx\svgscale\undefined%
      \relax%
    \else%
      \setlength{\unitlength}{\unitlength * \real{\svgscale}}%
    \fi%
  \else%
    \setlength{\unitlength}{\svgwidth}%
  \fi%
  \global\let\svgwidth\undefined%
  \global\let\svgscale\undefined%
  \makeatother%
  \parbox[c][][c]{\unitlength}{%
    \begin{picture}(1,0.6684636118598383)%
      \put(0,0){\includegraphics[width=\unitlength,viewport = 0 0 213 142]{feynman/lep_4.eps}}%
      \put(0.08301887357973989,0.06145553126168058){\raisebox{-0.5\height}{\rotatebox{0.0}{\scalebox{1.0}{\makebox(0,0)[]{{$e^+$}}}}}}%
      \put(0.0830188665107254,0.5903504050195699){\raisebox{-0.5\height}{\rotatebox{0.0}{\scalebox{1.0}{\makebox(0,0)[]{{$e^-$}}}}}}%
      \put(0.6639415915642487,0.16798471824888486){\raisebox{-0.5\height}{\rotatebox{0.0}{\scalebox{1.0}{\makebox(0,0)[]{{$\ell^+$}}}}}}%
      \put(0.8190705302569445,0.43929364096522094){\raisebox{-0.5\height}{\rotatebox{0.0}{\scalebox{1.0}{\makebox(0,0)[]{{$\ell^-$}}}}}}%
      \put(0.4743935309973046,0.40161725067385445){\raisebox{-0.5\height}{\rotatebox{0.0}{\scalebox{1.0}{\makebox(0,0)[]{{$\gamma^\ast$}}}}}}%
      \put(0.8776280436554367,0.564420488002808){\raisebox{-0.5\height}{\rotatebox{0.0}{\scalebox{1.0}{\makebox(0,0)[]{{$\gamma$}}}}}}%
      \put(0.9180512129380054,0.215633423180593){\raisebox{-0.5\height}{\rotatebox{0.0}{\scalebox{1.0}{\makebox(0,0)[]{{$\chi_{\ell J}$}}}}}}%
    \end{picture}%
  }%
\endgroup%

        \caption{}\label{fig:diag_leptonium_b}
    \end{subfigure}
    \caption{\it\small {Representative Feynman diagrams for leptonium production at $e^+e^-$ colliders: (a) para- or ortho-positronium production via $t$-channel photon exchange; (b) ortho-leptonium ($\ell=e,\mu,\tau$) production in association with a photon from final-state radiation. The diagrams were generated using {\tt FeynGame}~\cite{feyngame,feyngame2,feyngame3}.}}
    \label{fig:diag_leptonium}
\end{figure}
%%%%%%%%%%%%%%%%%%%%%%%%%%%%%%%%%%%%%%%%%%%%%%%%%%%%%%%%%%%%%%%%

With the upgraded implementation in {\madsons}, it is now possible to compute cross-section predictions not only for \swavetxt{} leptonia, as already demonstrated in ref.~\cite{ColpaniSerri:2025vdz}, but also for \pwavetxt{} configurations. To validate the NRQED implementation, we compare cross sections obtained with \mgshort{} against analytic calculations. In particular, we consider the processes $e^+e^- \to {\cal L}(2\pwave)+\gamma$ in the limit of infinitely large $Z$- and Higgs-boson masses, complementing the study of \swavetxt{} para- and ortho-positronium presented in figure~6 of ref.~\cite{ColpaniSerri:2025vdz}. The final-state leptonia ${\cal L}(2\pwave)$ considered here are para-positronium $h_e(2\pwave)$, corresponding to the ${}^1\pwave_1$ Fock state, as well as ortho-positronium $\chi_{eJ}(2\pwave)$ and ortho-ditauonium $\chi_{\tau J}(2\pwave)$, corresponding to the ${}^3\pwave_J$ Fock states. Para-ditauonium production, $e^+e^-\to h_\tau(2\pwave)+\gamma$, is instead forbidden by charge-conjugation symmetry in QED.
The representative Feynman diagrams contributing to these processes are shown in figure~\ref{fig:diag_leptonium}. All diagrams contribute to positronium production, while only the diagram in figure~\ref{fig:diag_leptonium_b} contributes to ditauonium production.

Before presenting the benchmark results, we comment on the different conventions for the principal quantum number used in quarkonium and leptonium spectroscopy. In quarkonium physics, the principal quantum number is conventionally defined as $N=n_r+1$, where $n_r$ denotes the radial quantum number, counting the number of radial nodes and vanishing for the ground state. In hydrogen-like Coulombic systems, however, the principal quantum number is defined such that the lowest \pwavetxt{} state corresponds to $N=2$. We adopt this latter convention for leptonia, as it naturally reflects the degeneracy of the Coulomb potential among states with the same principal quantum number, such as $N\swave$ and $N\pwave$ levels.

% ---------------- Positronium ----------------

Since the leptons forming the bound state must be kept massive, the \ufo{} model has to be loaded with the \texttt{lepton\_masses} restriction. In addition, the standard \mgshort{} syntax can be used to exclude diagrams involving $Z$ or Higgs bosons, thereby reproducing the limit in which electroweak bosons are infinitely heavy. For example, the para-positronium production process is generated via the commands
\prompt{import model sm\_onia-lepton\_masses}
\vskip -0.25truecm
\prompt{generate e+ e- > Ps(2|1P1) a / z h; output; launch}
\noindent
To account for the low energy scale relevant for leptonium production, we employ the $\alpha(0)$ scheme rather than the $G_\mu$ scheme, which is the default setup of the \texttt{sm\_onia} model. Since we only consider contributions involving internal quasi-real photons, this can be achieved by setting the first entry of \texttt{BLOCK SMINPUTS} in the \texttt{param\_card.dat} file to $137.036$, corresponding to the inverse electromagnetic coupling in the Thomson limit,
\begin{equation}
\alpha(0)=\dfrac{1}{137.036}\,.
\end{equation}
In all calculations, the leptonium mass is chosen as the sum of the constituent lepton masses,
\begin{equation}
M_{\mathcal L}=m_{\ell^-}+m_{\ell^{\prime+}}\,,
\end{equation}
%%%%%%%%%%%%%%%%%%%%%%%%%%%%%%%%%%%%%%%%%%%%%%%%%%%%%%%%%%%%%%%%
\vspace{\fill}
\begin{figure}[h]
    \centering
    \begin{subfigure}[b]{0.49\textwidth}
        \input{figures/xsec_paraEE.pgf}
        %\includegraphics{figures/xsec_paraEE.pdf}
        \caption{}
        \label{fig:xsec_EE_a}
    \end{subfigure}
    \begin{subfigure}[b]{0.49\textwidth}
        \input{figures/xsec_orthoEE0.pgf}
        %\includegraphics{figures/xsec_orthoEE0.pdf}
        \caption{}
        \label{fig:xsec_EE_b}
    \end{subfigure}
    \begin{subfigure}[b]{0.49\textwidth}
        \input{figures/xsec_orthoEE1.pgf}
        %\includegraphics{figures/xsec_orthoEE1.pdf}
        \caption{}
        \label{fig:xsec_EE_c}
    \end{subfigure}
    \begin{subfigure}[b]{0.49\textwidth}
        \input{figures/xsec_orthoEE2.pgf}
        %\includegraphics{figures/xsec_orthoEE2.pdf}
        \caption{}
        \label{fig:xsec_EE_d}
    \end{subfigure}
    \caption{\small\it Cross sections as functions of the c.m.\@ energy for the production in $e^+ e^-$ collisions of a photon in association with a \pwavetxt\ (a) para- and (b-d) ortho-positronium. In each figure, the upper panels show the total cross section obtained with {\tt\mgshort} (red dots) and from analytical expressions (blue line), whereas the lower panels display the ratio, normalised to the analytical reference values. Error bars on the {\tt \mgshort} points indicate statistical MC uncertainties.}
    \label{fig:xsec_EE}
\end{figure}
\vspace{\fill}
%%%%%%%%%%%%%%%%%%%%%%%%%%%%%%%%%%%%%%%%%%%%%%%%%%%%%%%%%%%%%%%%
\newpage\noindent
such that no momentum reshuffling is required. This choice also considerably simplifies the analytic evaluation of the cross section, allowing for a direct comparison with the \mgshort{} predictions. 

For $e^+e^-\to h_e(N\pwave)+\gamma$, the analytic cross section reads
\begin{equation}
\begin{aligned}
    & \lim_{m_Z,m_H\to\infty}\sigma\left(e^+e^-\to h_e(N\pwave)+\gamma\right)\\
    & \qquad = \dfrac{\pi\alpha^8}{24 s}\dfrac{N^2-1}{N^5}\dfrac{1}{\left(1 - \xi_{e}\right)^5}\, 
    \left\lbrace\sqrt{1 - \xi_{e}}\,\bigg\lbrack 36 - 110 \xi_{e} + 82 \xi_{e}^{2} + 18 \xi_{e}^{3} - 38 \xi_{e}^{4} \bigg\rbrack\right.\\
    &\qquad\quad \left. - \xi_{e}\log\!\left(\dfrac{1 - \sqrt{1 -\xi_{e}}}{1 + \sqrt{1 - \xi_{e}}}\right) \bigg\lbrack 16 - 5 \xi_{e} - 29 \xi_{e}^{2} + 31 \xi_{e}^{3} - 7 \xi_{e}^{4} \bigg\rbrack \right\rbrace\\
    & \quad 
    \overset{s\gg m_e^2}{\longrightarrow} \dfrac{3\pi\alpha^8}{2s}\dfrac{N^2-1}{N^5}\,,
\end{aligned}
\end{equation}
where $\xi_e=4m_e^2/s$.
The comparison of the total exclusive cross section is shown in figure~\ref{fig:xsec_EE_a} as a function of the ratio of the c.m.\@ energy to the positronium mass, $\sqrt{s}/2m_e$. The agreement between the analytic calculation and the independent \mgshort{} prediction is excellent over the entire energy range considered, from the production threshold ($\sqrt{s}\gtrsim 2m_e$) to the asymptotic high-energy regime ($\sqrt{s}\gg 2m_e$), where the electron mass becomes negligible. The relative differences remain well within the numerical MC uncertainties and are typically at the few-permille level.
Similar conclusions hold for the spin-triplet ortho-positronium production processes, $e^+e^-\to\chi_{eJ}+\gamma$, shown for $J=0,\ 1,\ \text{and}\ 2$ in figures~\ref{fig:xsec_EE_b}, \ref{fig:xsec_EE_c}, and \ref{fig:xsec_EE_d}, respectively. The analytic expressions used for comparison are given by
\begin{align}
    & \lim_{m_Z,m_H\to\infty}\sigma\left(e^+e^-\to\chi_{e 0}(N\pwave)+\gamma\right)\nonumber \\
    &\qquad = \dfrac{\pi\alpha^8}{216s}\dfrac{N^2-1}{N^5}\dfrac{1}{\left(1-\xi_e\right)^5}\left\lbrace\sqrt{1-\xi_e}\,\bigg\lbrack 236  - 550 \xi_{e} + 348\xi_{e}^{\!2} + 176 \xi_{e}^{3} - 406 \xi_{e}^{4} + 148 \xi_{e}^{5}\right.\nonumber \\
    & \qquad\quad\left. + 30 \xi_{e}^{6} - 18 \xi_{e}^{7} \bigg\rbrack +3\xi_{e}\log\!\left(\dfrac{1-\sqrt{1-\xi_e}}{1+\sqrt{1-\xi_e}}\right)\bigg\lbrack 4 - 23 \xi_{e} + 13 \xi_{e}^{2}+ 27 \xi_{e}^{3} - 45 \xi_{e}^{4} + 18 \xi_{e}^{5} \bigg\rbrack\right\rbrace\nonumber\\
    &\quad\overset{s\gg m_e^2}{\longrightarrow} \dfrac{59\pi\alpha^8}{54s}\dfrac{N^2-1}{N^5}\,,
\end{align}
\begin{align}
    & \lim_{m_Z,m_H\to\infty}\sigma\left(e^+e^-\to\chi_{e 1}(N\pwave)+\gamma\right)\nonumber\\
    &\qquad = \dfrac{\pi\alpha^8}{72s}\dfrac{N^2-1}{N^5}\dfrac{1}{\left(1-\xi_e\right)^5}\left\lbrace\sqrt{1-\xi_e}\,\bigg\lbrack 112 - 264 \xi_{e} + 304\xi_{e}^{2} - 206 \xi_{e}^{3} + 70 \xi_{e}^{4} \right.\nonumber\\
    &\qquad\quad\left. - 48 \xi_{e}^{5} - 4 \xi_{e}^{6} \bigg\rbrack + 3\xi_{e}\log\!\left(\dfrac{1-\sqrt{1-\xi_e}}{1+\sqrt{1-\xi_e}}\right)\bigg\lbrack 8 - 24 \xi_{e} + 36 \xi_{e}^{2} - 33 \xi_{e}^{3} + 7 \xi_{e}^{4} \bigg\rbrack\right\rbrace\nonumber\\
    &\qquad\overset{s\gg m_e^2}{\longrightarrow} \dfrac{14\pi\alpha^8}{9s}\dfrac{N^2-1}{N^5}\,,
\end{align}
\begin{align}
    &\lim_{m_Z,m_H\to\infty}\sigma\left(e^+e^-\to\chi_{e 2}(N\pwave)+\gamma\right)\nonumber\\
    &\qquad = \dfrac{\pi\alpha^8}{216s}\dfrac{N^2-1}{N^5}\dfrac{1}{\left(1-\xi_e\right)^5}\left\lbrace\sqrt{1-\xi_e}\,\bigg\lbrack 736 - 2528 \xi_{e} + 3912\xi_{e}^{2} - 1166 \xi_{e}^{3} - 1634 \xi_{e}^{4}\right.\nonumber\\
    &\qquad\quad + 512 \xi_{e}^{5} + 12 \xi_{e}^{6} - 24 \xi_{e}^{7}  \bigg\rbrack+3\xi_{e}\log\!\left(\dfrac{1-\sqrt{1-\xi_e}}{1+\sqrt{1-\xi_e}}\right)\bigg\lbrack 8 - 124 \xi_{e} + 608 \xi_{e}^{2} - 717 \xi_{e}^{3}\nonumber \\
    &\qquad\quad\left. + 195 \xi_{e}^{4} \bigg\rbrack\right\rbrace \nonumber\\
    &\quad\overset{s\gg m_e^2}{\longrightarrow} \dfrac{92\pi\alpha^8}{27s}\dfrac{N^2-1}{N^5}\,.
\end{align}
In all cases, the \mgshort{} predictions agree with the analytic results within statistical MC uncertainties over the full energy range considered. This provides a strong validation of the implementation of the \pwavetxt{} projectors and LDMEs for NRQED bound states.

% ---------------- Ditauonium ----------------
%%%%%%%%%%%%%%%%%%%%%%%%%%%%%%%%%%%%%%%%%%%%%%%%%%%%%%%%%%%%%%%%
\begin{figure}[t]
    \centering
    \begin{subfigure}[b]{0.49\textwidth}
        \input{figures/xsec_orthoTT0.pgf}
        %\includegraphics{figures/xsec_orthoTT0.pdf}
        \caption{}
        \label{fig:xsec_TT_a}
    \end{subfigure}
    \begin{subfigure}[b]{0.49\textwidth}
        \input{figures/xsec_orthoTT1.pgf}
        %\includegraphics{figures/xsec_orthoTT1.pdf}
        \caption{}
        \label{fig:xsec_TT_b}
    \end{subfigure}
    \begin{subfigure}[b]{0.49\textwidth}
        \input{figures/xsec_orthoTT2.pgf}
        %\includegraphics{figures/xsec_orthoTT2.pdf}
        \caption{}
        \label{fig:xsec_TT_c}
    \end{subfigure}
    \caption{\small\it Same as figure~\ref{fig:xsec_EE}, but for $\chi_{\tau J}$ associated production. The state $h_{\tau}(2\mathrm{P})$ is not considered, as the process $e^+e^-\to h_{\tau}(2\mathrm{P})+\gamma$ is forbidden in QED by charge-conjugation symmetry.}
    \label{fig:xsec_TT}
\end{figure}
%%%%%%%%%%%%%%%%%%%%%%%%%%%%%%%%%%%%%%%%%%%%%%%%%%%%%%%%%%%%%%%%

Owing to the absence of the $t$-channel diagrams shown in figure~\ref{fig:diag_leptonium_a}, the analytic expressions for the ortho-ditauonium production processes $e^+e^-\to\chi_{\tau J}(N\pwave)+\gamma$ take a particularly compact form,
\begin{equation}
    \lim_{m_Z,m_H\to\infty}\sigma\left(e^+e^-\to\chi_{\tau 0}(N\pwave)+\gamma\right) = \dfrac{\pi\alpha^8}{54s}\dfrac{N^2-1}{N^5}\dfrac{\xi_\tau\left(1-3\xi_\tau\right)^2}{(1-\xi_\tau)}\,,
\end{equation}
\begin{equation}
    \lim_{m_Z,m_H\to\infty}\sigma\left(e^+e^-\to\chi_{\tau 1}(N\pwave)+\gamma\right) = \dfrac{\pi\alpha^8}{9s}\dfrac{N^2-1}{N^5}\dfrac{\xi_\tau\left(1+\xi_\tau\right)}{(1-\xi_\tau)}\,,
\end{equation}
\begin{equation}
    \lim_{m_Z,m_H\to\infty}\sigma\left(e^+e^-\to\chi_{\tau 2}(N\pwave)+\gamma\right) = \dfrac{\pi\alpha^8}{27s}\dfrac{N^2-1}{N^5}\dfrac{\xi_\tau\left(1+3 \xi_\tau+6\xi_\tau^2\right)}{(1-\xi_\tau)}\,,
\end{equation}
where $\xi_\tau = 4m_\tau^2/s$. The comparison between the analytic results and the \mgshort{} predictions for the three total angular momentum states $J=0,\ 1,\ \text{and}\ 2$ is shown in figures~\ref{fig:xsec_TT_a}, \ref{fig:xsec_TT_b}, and \ref{fig:xsec_TT_c}, respectively.
Over the entire kinematic range, the analytic predictions are reproduced by \mgshort{} at the sub-permille level, fully consistent with the quoted MC uncertainties. This provides further strong evidence for the validity of the implementation.

Before concluding this section, we elaborate further on the $\chi_{\tau J}$ results.
While the cross sections for $\chi_{\tau 1}(2\pwave)$ and $\chi_{\tau 2}(2\pwave)$ exhibit a smooth behaviour over the full energy range, the $\chi_{\tau 0}(2\pwave)$ production cross section vanishes at $\sqrt{s}=\sqrt{3}\,M_{\chi_{\tau 0}(2\pwave)}$. This feature is related to the composite nature of the bound state and is also present in the analogous charmonium and bottomonium production processes~\cite{Chung:2008km,Li:2009ki,Sang:2009jc,Xu:2014zra,Brambilla:2017kgw}. At LO, the observed cancellation arises from a dynamical mechanism in which different terms cancel at the amplitude level. Higher-order corrections in both the strong/electromagnetic coupling and velocity expansions are not expected to preserve this exact cancellation, so improved predictions would no longer yield a vanishing cross section at the same c.m.\@ energy. Nevertheless, since these are only subleading corrections, a pronounced dip is expected to survive, providing a characteristic signature that could be studied to test \pwavetxt\ leptonium and charmonium production experimentally, for example at future Super Tau-Charm Facilities~\cite{Achasov:2023gey,SCTF_CDR_PhysicsDetector_2018}. By contrast, this cancellation is not a generic feature of scalar particle production in association with a photon. For instance, no analogous dip occurs in the process $e^+e^-\to H+\gamma$ in an effective field theory containing an effective $\gamma\gamma H$ interaction. 